problem:  
Let $\Sigma^n$ be a smooth manifold homeomorphic to $S^n$ with $n\ge7$ (hence a homotopy sphere). Suppose $\Sigma^n$ admits a Riemannian metric $g$ with strictly positive sectional curvature. The Differentiable Sphere Theorem (Brendle–Schoen) asserts that if $g$ is $1/4$‑pinched, then $\Sigma^n$ is diffeomorphic to the standard sphere.  

Now fix $\lambda\in(0,1/4)$ and let $\mathcal{P}_\lambda$ be the class of Riemannian metrics with  
\[
\frac{K_{\min}(p)}{K_{\max}(p)}\ge \lambda>0
\quad \text{for all points } p.
\]
Assume that $\Sigma^n$ possesses some $g\in\mathcal{P}_\lambda$.  
Consider the Ricci flow $\partial_t g_t = -2\,\mathrm{Ric}(g_t)$ with $g_0 = g$.

**Question:**  
Assuming short‑time existence for $g_t$, does the Ricci flow improve the pinching ratio uniformly?  
More precisely, is there $\varepsilon(\lambda)>0$ such that any $\lambda$‑pinched positively curved metric on a homotopy sphere evolves under Ricci flow to become $(\lambda+\varepsilon(\lambda))$‑pinched before forming any singularity?  
If such a curvature improvement always occurs, can one deduce that any positively curved exotic sphere would flow (before finite time singularity) into a regime covered by the Differentiable Sphere Theorem, thereby forcing it to have the standard smooth structure?  

Why is it a "good" problem:  
This problem is conceptually clean yet touches the frontiers of geometric analysis and topology. It transforms the longstanding question of positive curvature on exotic spheres into a **dynamical pinching‑improvement problem under Ricci flow**—a quantitative, analytic reformulation rather than a purely existence‑based one. The question is open and subtle: even short‑time curvature evolution on homotopy spheres under Ricci flow is poorly understood in the absence of strong symmetries.  

A positive answer would yield a new route toward classifying smooth structures on spheres via curvature evolution, potentially bridging Ricci flow analysis and high‑dimensional differential topology in a striking way. The statement is simple, logically precise, and mathematically sound, yet its resolution would represent major progress on fundamental geometric classification problems—thereby fulfilling the hallmarks of a “good” research problem.